\section{核心风险}

\begin{enumerate}
  \item \textbf{资金口径风险}：所谓主力资金是成交单大小分类，无法穿透到机构、量化、产业资本或个人账户。连续流入只能说明交易结构，不能证明长期持有人建仓。
  \item \textbf{低位钝化风险}：低位可以维持很久。银行若缺乏增量催化，估值修复可能只通过股息实现。
  \item \textbf{银行资产质量风险}：区域地产、零售经营贷和消费信用风险可能使不良率、关注率和信用成本恶化。
  \item \textbf{息差风险}：贷款重定价可能快于存款成本下降，当前息差改善未必持续。
  \item \textbf{徐工周期风险}：国内工程机械需求、出口竞争、海外贸易政策和矿山资本开支均可能反转。
  \item \textbf{汇率风险}：徐工海外占比高，人民币升值会影响汇兑和利润折算。
  \item \textbf{风格风险}：如果高位科技重新进入全面主升，低位红利和周期修复可能继续缺乏关注。
  \item \textbf{已启动拥挤风险}：中兴和江波龙近期成交与波动明显放大，即便长期逻辑正确，也可能先经历20\%以上回撤。
  \item \textbf{盈利桥风险}：江波龙H1预告未经审计，中兴算力业务毛利率未披露，恒瑞管线概率和未来BD确认时点存在不确定性。
  \item \textbf{工业富联兑现风险}：H1预告未经审计，AI服务器客户集中、GPU平台量产、低毛利制造属性和营运资金可能使H2利润低于预期。
  \item \textbf{旧分母失效风险}：香农芯创、兆易创新、多氟多和方正科技H1预告EPS已超过旧全年EPS，继续引用旧PE会产生方向性错误。
\end{enumerate}

\section{监控仪表盘}

\begin{exhibitbox}[Exhibit 24：下一季度验证与降级规则]
\begin{tabularx}{\textwidth}{L{2.2cm}L{3.8cm}L{3.8cm}X}
\toprule
\textbf{标的/板块} & \textbf{升级信号} & \textbf{降级信号} & \textbf{动作} \\
\midrule
渝农商行 & 营收增速>4\%，净息差$\geq$1.60\%，不良率$\leq$1.10\% & 净息差<1.55\%或不良率>1.15\% & 升级/降为仅收息观察 \\
徐工机械 & 海外增速>10\%，现金流持续为正，调整后利润恢复双位数 & EPS下修至0.60以下，海外增速<8\% & 加仓确认/退出增长估值 \\
沪农商行 & 营收继续正增长，拨备覆盖率$\geq$300\%，分红率稳定 & 营收转负或拨备<280\% & 收益配置/降为纯观察 \\
银行板块 & 周度资金连续流入、银行ETF持续净申购 & 周度资金转负且长端利率/信用风险恶化 & 保持/降低行业权重 \\
军工/传媒 & 回踩不破启动位、资金由事件扩散至业绩标的 & 一日游、ETF份额回吐 & 右侧跟踪/不追高 \\
半导体设备 & 估值消化、价格回落至中位、订单继续兑现 & 高位放量流出、盈利预测下修 & 等待/回避 \\
中兴通讯 & 算力占比$\geq$25\%、H2利润恢复、现金流改善 & EPS<1.30、现金流持续为负 & 回踩确认/降为期权观察 \\
江波龙 & H2净利$\geq$80亿元、现金流改善、无大额减值 & 全年净利<160亿元或库存/价格反转 & 高弹性配置/取消增长信用 \\
恒瑞医药 & 创新药销售$\geq$30\%、EPS$\geq$1.35、里程碑兑现 & 创新药增速<20\%或管线延期 & 趋势配置/下修SOTP \\
工业富联 & H2净利$\geq$321亿元、AI服务器增速>100\%、毛利率$\geq$7\% & 平台量产延期、净利<516亿元或毛利率<7\% & 回撤配置/下修PE \\
预告重估池 & 新研报采用扣非与新股本分母，现金流同步改善 & 继续沿用预告前旧EPS或一次性收益主导 & 升级建模/维持观察 \\
\bottomrule
\end{tabularx}
\end{exhibitbox}

\section{最终结论}

当前市场不能只看低位，也不能只追强势。本报告用两条曲线覆盖机会：渝农商行、徐工机械、沪农商行代表尚未充分启动的安全边际；中兴通讯、江波龙、恒瑞医药和工业富联代表趋势已启动但盈利分母仍能扩张的机会。紫光股份、香农芯创、兆易创新、方正科技和中船防务进入预告后重估池，不在分母刷新前伪造目标。军工、先进封装和传媒已经启动，但尚未同时通过公司利润桥与现价空间门槛。真正的执行纪律是\textbf{静默篮子左侧分批，启动篮子只做回撤和业绩确认，预告重估池只等新分母}。
